% =============================================================================
% Compact template for integration with an empty agent: Nanjing University of Science and Technology bachelor's thesis format
% Cover all structural components with minimal body text, focusing on format reuse
% Compile: tectonic main.tex or xelatex && bibtex && xelatex && xelatex
% =============================================================================
\documentclass[a4paper, 12pt, twoside]{ctexart}

% --- Page setup (A4; top 30mm, bottom 24mm, left/right 25mm; header footskip=20mm) ---
\usepackage[
  a4paper,
  top=30mm, bottom=24mm, left=25mm, right=25mm,
  headheight=15pt, headsep=10mm, footskip=20mm
]{geometry}

% --- Body: SimSun, Chinese small size 4, approximately 20pt fixed line spacing ---
\usepackage{setspace}
\AtBeginDocument{\zihao{-4}\songti\setlength{\baselineskip}{20pt}\setlength{\parindent}{2em}}

% --- Prevent Overfull \hbox by allowing paragraph word spacing to stretch ---
\emergencystretch=0.5em

% --- Font settings ---
% Chinese: SimSun for body text; synthetic bold SimSun for headings to meet the school's bold-SimSun requirement
% English and numbers: Times New Roman
\usepackage{fontspec}
\setmainfont{Times New Roman}
\setCJKmainfont[AutoFakeBold=2.5]{Songti SC}
\setCJKfamilyfont{song}[AutoFakeBold=2.5]{Songti SC}
\setCJKfamilyfont{hei}{Heiti SC}

% --- Mathematical formulas ---
\usepackage{amsmath, amssymb}

% --- Figures and tables ---
\usepackage{graphicx}
\graphicspath{{../图表/}}
\usepackage{float}
\usepackage{subcaption}
\DeclareCaptionFont{songwu}{\songti\zihao{5}}
\captionsetup{font=songwu, labelfont=songwu, labelsep=space}
\captionsetup[lstlisting]{font=songwu, labelfont=songwu, labelsep=space}
\usepackage{booktabs}
\usepackage{multirow}
\usepackage{longtable}
\usepackage{tabularx}
\usepackage{adjustbox}
\usepackage{afterpage}

% --- Number figures, tables, and equations by chapter (Figure X-X, Table X-X, Equation X-X) ---
\IfFileExists{chngcntr.sty}{\usepackage{chngcntr}}{}
\counterwithin{figure}{section}
\counterwithin{table}{section}
\counterwithin{equation}{section}
\def\thefigure{\thesection.\arabic{figure}}
\def\thetable{\thesection.\arabic{table}}
\def\theequation{\thesection.\arabic{equation}}

% --- Code blocks ---
\usepackage{listings}
\usepackage{listingsutf8}
\usepackage{xcolor}
\lstset{
    basicstyle=\zihao{5}\ttfamily,
    columns=fullflexible,
    keepspaces=true,
    literate={，}{{，}}1 {。}{{。}}1 {：}{{：}}1 {；}{{；}}1
            {（}{{（}}1 {）}{{）}}1 {、}{{、}}1,
    keywordstyle=\color{blue},
    commentstyle=\color{gray},
    stringstyle=\color{red!70!black},
    numbers=left,
    numberstyle=\tiny\color{gray},
    frame=single,
    breaklines=true,
    tabsize=4,
}

% --- Cross-references and hyperlinks (load hyperref as late as possible) ---
\usepackage{hyperref}
\hypersetup{
    colorlinks=true,
    linkcolor=black,
    citecolor=black,
    urlcolor=black,
}

% --- References (GB/T 7714-2015 numeric style; sort&compress compacts citation numbers) ---
% Use natbib + gbt7714-numerical; for BibTeX, run xelatex twice, bibtex, then xelatex twice
\usepackage[numbers,sort&compress]{natbib}
\bibliographystyle{gbt7714-numerical}

% --- Headers and footers (different odd/even pages, page numbers on the outside) ---
% Even pages (CE): chapter title on the left, "Bachelor's Thesis" on the right
% Odd pages (CO): "Bachelor's Thesis" on the left, chapter title on the right
% Page numbers at LE (left even) and RO (right odd), i.e. the outer edges
\usepackage{fancyhdr}
\pagestyle{fancy}
\fancyhf{}
% Correct sectionmark definition; must precede its use
\makeatletter
\def\sectionmark#1{\markboth{\thesection\quad #1}{}}
\makeatother
\fancyhead[CO]{\songti\zihao{5} 学士学位论文\hfill <论文题目(简)>}
\fancyhead[CE]{\songti\zihao{5} \leftmark\hfill 学士学位论文}
\fancyfoot[LE,RO]{\songti\zihao{5}\thepage}
\fancypagestyle{plain}{%
  \fancyhf{}%
  \fancyhead[CO]{\songti\zihao{5} 学士学位论文\hfill <论文题目(简)>}%
  \fancyhead[CE]{\songti\zihao{5} \leftmark\hfill 学士学位论文}%
  \fancyfoot[LE,RO]{\songti\zihao{5}\thepage}%
}

% --- TOC depth: include three heading levels ---
% tocloft: \cftsecnumdepth=3 includes section/sub/subsub levels
\IfFileExists{tocloft.sty}{\usepackage[titles]{tocloft}}{}
\def\contentsname{目\quad 录}
\def\refname{参考文献}

% --- Chapter heading format (school requirements) ---
\ctexset{
    section = {
        format = \raggedright\songti\bfseries\zihao{-3},
        beforeskip = 18pt,
        afterskip = 18pt,
        break = \clearpage,
    },
    subsection = {
        format = \raggedright\songti\bfseries\zihao{4},
        beforeskip = 12pt,
        afterskip = 12pt,
    },
    subsubsection = {
        format = \raggedright\songti\bfseries\zihao{-4},
        beforeskip = 6pt,
        afterskip = 6pt,
    },
    paragraph = {
        format = \raggedright\songti\zihao{-4},
        beforeskip = 6pt,
        afterskip = 0pt,
    },
}

% --- Custom commands ---
\newcommand{\frontmattertitle}[1]{%
  \vspace*{18pt}\begin{center}{\songti\bfseries\zihao{3}#1}\end{center}\vspace{18pt}}
\newcommand{\centeredsectionstar}[1]{%
  \clearpage\phantomsection\vspace*{18pt}\begin{center}{\songti\bfseries\zihao{3}#1}\end{center}\vspace{18pt}}

% =============================================================================
\begin{document}

% ========================== Abstract ==========================
\newpage
\pagestyle{plain}
\pagenumbering{roman}

\frontmattertitle{摘\quad 要}
\phantomsection
\addcontentsline{toc}{section}{摘要}

本文针对XXXXXXXXXXXXXXXXXXXXXXXXXXXXXXXXXXXXXXXXXXXXXXXXXXXXXXXXXXXXXXXXXXXXXXXXXX。
XXXXXXXXXXXXXXXXXXXXXXXXXXXXXXXXXXXXXXXXXXXXXXXXXXXXXXXXXXXXXXXXXXXXXXXXXXXXXXXXXXX。
XXXXXXXXXXXXXXXXXXXXXXXXXXXXXXXXXXXXXXXXXXXXXXXXXXXXXXXXXXXXXXXXXXXXXXXXXXXXXXXXXXX。

\vspace{1cm}
\noindent{\songti\bfseries\zihao{4} 关键词：}{\songti\zihao{-4} 关键词1；关键词2；关键词3}

% ========================== Abstract ==========================
\newpage
\markboth{Abstract}{}
\begin{center}{\bfseries\zihao{3} Abstract}\end{center}\vspace{18pt}
\phantomsection
\addcontentsline{toc}{section}{Abstract}

This thesis presents a systematic study of XXXXXXXXXXXXXXXXXXXXXXXXXXXXXXXXXXXXXXXX.
XXXXXXXXXXXXXXXXXXXXXXXXXXXXXXXXXXXXXXXXXXXXXXXXXXXXXXXXXXXXXXXXXXXXXXXXXXXXXXXXXXX.
XXXXXXXXXXXXXXXXXXXXXXXXXXXXXXXXXXXXXXXXXXXXXXXXXXXXXXXXXXXXXXXXXXXXXXXXXXXXXXXXXXX.

\vspace{1cm}
\noindent{\bfseries\zihao{4} Keywords:}{\zihao{-4} keyword1; keyword2; keyword3}

% ========================== Table of contents ==========================
\newpage
\markboth{目\quad 录}{}
\begingroup
\def\contentsname{目\quad 录}
\tableofcontents
\endgroup
% For duplex printing, add a blank page if the front matter has an odd page count
\ifodd\value{page}\newpage\thispagestyle{empty}\mbox{}\fi

% ========================== Main text ==========================
\newpage
\pagestyle{fancy}
\pagenumbering{arabic}

% ========================================================================
%   Chapter 1: Introduction
% ========================================================================
\section{绪论}

\subsection{研究背景与意义}

XXXXXXXXXXXXXXXXXXXXXXXXXXXXXXXXXXXXXXXXXXXXXXXXXXXXXXXXXXXXXXXXXXXXXXXXXXXXX。
如图\ref{fig:architecture}所示，XXXXXXXXXXXXXXXXXXXXXXXXXXXXXXXXXXXXXXXXX。

\begin{figure}[H]
    \centering
    \includegraphics[width=0.65\textwidth]{architecture.jpg}
    \caption{系统架构图}
    \label{fig:architecture}
\end{figure}

\subsection{国内外研究现状}

XXXXXXXXXXXXXXXXXXXXXXXXXXXXXXXXXXXXXXXXXXXXXXXXXXXXXXXXXXXXXXXXXXXXXXXXXXXXX。
文献\cite{author2024example}提出了XXXXXXXXXXXXXXXXXXXXXXXXXXXXXXXXXXXXXXXXXX。

\subsubsection{Linux调度器研究}

XXXXXXXXXXXXXXXXXXXXXXXXXXXXXXXXXXXXXXXXXXXXXXXXXXXXXXXXXXXXXXXXXXXXXXXXXXXXX。

% Use paragraph for level-four headings; school requirements generally omit subsubsubsection
\paragraph{相关技术基础}

XXXXXXXXXXXXXXXXXXXXXXXXXXXXXXXXXXXXXXXXXXXXXXXXXXXXXXXXXXXXXXXXXXXXXXXXXXXX。

\subsection{本文创新点}

本文主要创新点如下：

\begin{enumerate}
    \item 创新点一：XXXXXXXXXXXXXXXXXXXXXXXXXXXXXXXXXXXXXXXXXXXXXXXX。
    \item 创新点二：XXXXXXXXXXXXXXXXXXXXXXXXXXXXXXXXXXXXXXXXXXXXXXXX。
\end{enumerate}

% ========================================================================
%   Chapter 2: Technical Background
% ========================================================================
\section{相关技术基础}

本章介绍XXXXXXXXXXXXXXXXXXXXXXXXXXXXXXXXXXXXXXXXXXXXXXXXXXXXXXXXXXXXXXXX。

\subsection{Linux调度器原理}

CFS的核心思想是XXXXXXXXXXXXXXXXXXXXXXXXXXXXXXXXXXXXXXXXXXXXXXXXXXXXXXXXX。
虚拟运行时间的计算如式\eqref{eq:vruntime}所示。

\begin{equation}
    \Delta vruntime = \Delta t_{wall} \times \frac{NICE\_0\_LOAD}{weight(task)}
    \label{eq:vruntime}
\end{equation}

表\ref{tab:cfs-params}列出了CFS的核心可调参数及其语义。

\begin{table}[H]
    \centering
    \caption{CFS核心可调参数}
    \label{tab:cfs-params}
    \begin{tabular}{lll}
        \toprule
        \textbf{参数名称} & \textbf{基准值} & \textbf{语义} \\
        \midrule
        \texttt{sched\_latency\_ns} & 6\,ms & 调度周期目标 \\
        \texttt{sched\_min\_granularity\_ns} & 0.75\,ms & 最小抢占粒度 \\
        \bottomrule
    \end{tabular}
\end{table}

% ========================================================================
%   Chapter 3: System Design
% ========================================================================
\section{系统设计}

XXXXXXXXXXXXXXXXXXXXXXXXXXXXXXXXXXXXXXXXXXXXXXXXXXXXXXXXXXXXXXXXXXXXXXXXXXXXX。

\subsection{总体架构}

XXXXXXXXXXXXXXXXXXXXXXXXXXXXXXXXXXXXXXXXXXXXXXXXXXXXXXXXXXXXXXXXXXXXXXXXXXXXX。

\section{实验与分析}

XXXXXXXXXXXXXXXXXXXXXXXXXXXXXXXXXXXXXXXXXXXXXXXXXXXXXXXXXXXXXXXXXXXXXXXXXXXXX。

% ========================================================================
%   References
% ========================================================================
% Reference page numbers continue in Arabic numerals from the main text; do not switch back to Roman numerals
% For BibTeX, replace \begin{thebibliography}... with
%   \bibliography{refs}, then run xelatex && bibtex && xelatex && xelatex
\newpage
\begin{thebibliography}{99}
  \bibitem{author2024example} Author A, Author B. Title of the reference[J].
    Journal Name, 2024, 1(1): 1-10. DOI: 10.0000/example.
  \bibitem{linuxkernel} Linux Kernel Documentation[EB/OL].
    \url{https://www.kernel.org/doc/html/latest/}.
\end{thebibliography}

% ========================== Appendix ==========================
\appendix
\centeredsectionstar{附\quad 录}
\markboth{附\quad 录}{}
\addcontentsline{toc}{section}{附录}

\subsection*{附录A：实验配置}

表~\ref{tab:appendix-configs}列出了实验配置。

\begin{table}[H]
    \centering
    \caption{实验配置表}
    \label{tab:appendix-configs}
    \begin{tabular}{llr}
        \toprule
        \textbf{参数} & \textbf{符号} & \textbf{值} \\
        \midrule
        CPU核心数 & $N$ & 4 \\
        调度周期 & $T$ & 10\,ms \\
        \bottomrule
    \end{tabular}
\end{table}

\subsection*{附录B：核心代码}

\begin{lstlisting}[language=C, caption={示例代码}, label={lst:sample}]
#include <stdio.h>
int main() {
    printf("Hello, World!\n");
    return 0;
}
\end{lstlisting}

\end{document}
